% Discussion for item2_2 -- prose goes here.
% scripts/make_latex.py will NEVER overwrite this file.

\textsc{ctrl+z} sends \texttt{SIGTSTP} to the foreground job and the shell
reports \texttt{[1]+ Stopped}. The process has \emph{not} been killed: \texttt{ps}
shows state \texttt{T}, stopped by a job-control signal. It is still in the
process table, still holds its memory and open files, and still has its place in
the job list --- it simply will not be scheduled again until told to continue,
which is why \textsc{ctrl+z} is not a way to abandon a program.

\texttt{fg} and \texttt{bg} both resume it by sending \texttt{SIGCONT}, and
differ only in who gets the terminal afterwards. \texttt{bg} resumes it in the
background: the state goes from \texttt{T} back to \texttt{S}, \texttt{jobs}
reports \texttt{Running}, and the prompt stays ours --- as though it had been
started with \texttt{\&} in the first place. (\texttt{S} is \emph{sleeping}, not
running on a CPU: the script is blocked in \texttt{sleep(1000)}. What changed is
that it is back under the scheduler rather than held off it.) \texttt{fg}
instead hands the terminal back and blocks until the job ends, exactly like the
first run of situation one --- so the pair is what makes \textsc{ctrl+z} useful
rather than merely disruptive: suspend, move to the background, reclaim the
terminal, without restarting the work.
